\section{Mở đầu}
\subsection{Lý do chọn đề tài}
\begin{itemize}
    \item Thống kê của Tổ chức Ung thư toàn cầu (GLOBOCAN) cho thấy, mỗi năm nước ta ghi nhận khoảng 21.555 ca mắc mới ung thư vú, chiếm \(25\% \) tổng số các bệnh ung thư ở nữ giới.
    Ở Việt Nam, ung thư vú cũng là loại ung thư đứng hàng đầu, chiếm tới 25,8\% các trường hợp ung thư ở nữ giới với 21.555 người mới được phát hiện bệnh và 9.345 bệnh nhân tử vong. 
    \item Hiện nay, ở Việt Nam, đội ngũ bác sĩ lành nghề chưa đáp ứng đủ nhu cầu xã hội, nhất là ở tuyến Huyện và Tỉnh. Việc thiếu đội ngũ Y bác sĩ chất lượng cao dẫn tới việc chẩn đoán các bệnh về ung thư còn chưa có độ chính xác cao (ở các tuyến huyện và tuyến dưới hơn). Việc chẩn đoán sai bệnh có thể gây ra nhiều hậu quả đáng tiếc
    \item Việc chẩn đoán đúng nguy ung thư lành tính và ác tính có ý nghĩa vô cùng quan trọng trong việc hỗ trợ các bác sĩ, chuyên gia y tế và bệnh nhân trong việc phát hiện và can thiệp kịp thời, giảm thiểu nguy cơ tử vong cũng như giảm bớt hậu quả lâu dài do bệnh ung thư gây ra.
\end{itemize}



\subsection{Mục tiêu dự án}
Nhận thấy sự cấp thiết đó, nhóm em đã tiến hành xây dựng mô hình chẩn đoán  ung thư vú(lành tính và ác tính )  sử dụng một số phương pháp học máy cơ bản. Mô hình này nhằm hỗ trợ phát hiện sớm các nguy cơ ung thư vú, giúp các bác sĩ và bệnh nhân có thể can thiệp kịp thời, giảm thiểu rủi ro tử vong cũng như hạn chế các di chứng nghiêm trọng có thể xảy ra. Với các phương pháp học máy: Naive Bayes, Random Forest, và Logistic Regression, mô hình được kỳ vọng sẽ cung cấp công cụ dự đoán chính xác, dễ triển khai và có thể ứng dụng thực tiễn trong môi trường y tế.
